\documentclass[12pt,UTF8]{ctexart}

\usepackage[a4paper,margin=1in]{geometry}
\usepackage{amsmath,amssymb,amsthm,mathtools,bm}
\usepackage{booktabs,tabularx,array,longtable}
\usepackage{enumitem}
\usepackage{setspace}
\usepackage{hyperref}

\hypersetup{hidelinks}
\allowdisplaybreaks[3]
\setlength{\emergencystretch}{3em}
\setlength{\extrarowheight}{2pt}
\setstretch{1.18}
\sloppy

\newcommand{\E}{\mathbb{E}}
\newcommand{\Prb}{\mathbb{P}}
\newcommand{\dd}{\,\mathrm{d}}
\newcommand{\ind}{\mathbf{1}}
\newcommand{\Rset}{\mathcal{R}}

\theoremstyle{plain}
\newtheorem{proposition}{命题}[section]
\newtheorem{assumption}{假设}[section]

\theoremstyle{remark}
\newtheorem{remark}{说明}[section]

\title{MAH 模型详细讲义\\
\large 路线扩张、实施强度、CMO 支持市场与观测份额}
\author{工作区：D:\textbackslash working\_paper\textbackslash DANNI}
\date{2026年7月17日}

\begin{document}

\maketitle
\tableofcontents
\bigskip

\section{导言：这份讲义到底要解释什么}

这份讲义解释的不是一个“所有市场都联立清算”的大模型，而是一个更聚焦、
也更适合当前论文问题的机制模型。核心问题只有一句话：

\begin{quote}
当 MAH 制度允许“持证主体”和“生产主体”分离以后，原研项目的商业化路线
会不会变多、变顺，从而提高企业今天做原研研发的私人回报？
\end{quote}

这里的经济逻辑有三层。

\begin{enumerate}[leftmargin=2em]
    \item 第一层是\textbf{离散的路线扩张}。MAH 以前，很多企业如果没有足够的
    自有生产能力，就只能内部生产、转让，或者放弃。MAH 以后，合法的
    “持证但委托生产”的 retained external route 出现了，记为路线 $E$。
    \item 第二层是\textbf{连续的实施强度}。即使法律上开放了 $E$ 路线，
    各地区、年份、剂型、监管执行环境也未必一样顺畅。我们用
    $\eta\in[0,1]$ 表示 post-MAH implementation intensity。$\eta$ 越高，
    说明在 MAH 规则已经存在的前提下，外部委托这条 retained route
    越容易真正跑通。
    \item 第三层是\textbf{CMO 支持市场的价格反馈}。如果很多企业都想走 $E$
    路线，对合格委托生产支持的需求会上升，这会推高有效支持成本 $p^*$，
    进而部分抵消制度改革的好处。这就是 attenuation。
\end{enumerate}

因此，本讲义的基准视角是：

\begin{quote}
MAH 不是一个一般性 demand shock，不是一个一般性 scientific productivity
shock，也不是一个纯粹的 approval-delay reform。它首先改变 route set，
然后通过 route-specific implementation 和 friction 影响原研研发激励。
\end{quote}

下面我们会把这个逻辑拆成本科生也能完整跟上的几个模块：制度背景、
对象与假设、收益函数、EV1 logit 推导、R\&D 最优化、CMO 市场闭合、
比较静态、观测份额、Bellman 递推、校准与识别、机制区分以及局限。

\paragraph{本节小结。}
这份模型讲义只研究一条窄机制：MAH 如何通过 retained external route
的可行性与顺畅度，影响原研项目的预期商业化价值和研发投入。

\section{制度问题：为什么 MAH 会进入原研研发激励模型}

\subsection{制度直觉}

在药品制度环境里，很多研究导向企业的问题不是“想不想创新”，而是
“即使项目做出来，最后能不能以自己保留 holder 身份的方式，把产品推到市场”。
如果必须自己生产，那么没有完备生产能力的研究导向企业就会在研发一开始
就预期到商业化受限。这样，当前的研发回报会被压低。

MAH 的制度意义正在这里：它把“谁持有上市许可”和“谁负责实际生产”
从严格绑定，变成了可以在监管规则下分离。经济上，这不是一句简单的
“允许外包”，因为 holder 责任并没有消失。持证方仍然要承担质量、
持续监督、技术转移、协议管理等责任。因此，外部路线不是无摩擦的。
这也是为什么模型里必须同时保留：

\begin{enumerate}[leftmargin=2em]
    \item 一项路线层面的\textbf{有效支持成本} $p$；
    \item 一项 MAH 可改善、但不一定完全消失的
    \textbf{路线摩擦} $\tau_i^E(M,\eta)$；
    \item 一项 holder 仍然承担的\textbf{剩余负担} $\mu_i^E$。
\end{enumerate}

\subsection{为什么要区分 \texorpdfstring{$M$ 和 $\eta$}{M 和 eta}}

本讲义把改革分成两个对象。

\begin{enumerate}[leftmargin=2em]
    \item $M\in\{0,1\}$：制度是否在法律上开放 MAH retained external route。
    $M=0$ 是 pre-MAH，$M=1$ 是 post-MAH。
    \item $\eta\in[0,1]$：在 $M=1$ 条件下，实施强度有多高。$\eta$ 可以理解成
    地区、年份、剂型、监管执行、CMO 配套成熟度等综合指标。
\end{enumerate}

这样做的好处是：

\begin{enumerate}[leftmargin=2em]
    \item $M$ 专门处理\textbf{离散的路线开放}；
    \item $\eta$ 专门处理\textbf{post-MAH 内部的连续差异}；
    \item 经验上也更贴近数据：很多时候 pre-MAH 根本没有可比的 $E$ 路线，
    但 post-MAH 不同地区之间存在不同的 implementation intensity。
\end{enumerate}

\paragraph{本节小结。}
MAH 在模型里不是一般性“利好政策”，而是一个改变商业化路线可行性和
外部生产实施条件的制度安排，因此要同时区分法律开放和实施强度。

\section{模型对象与核心假设}

\subsection{对象表}

\begingroup
\small
\renewcommand{\arraystretch}{1.14}
\noindent\begin{tabularx}{\textwidth}{p{0.13\textwidth} p{0.17\textwidth} X p{0.14\textwidth}}
\toprule
符号 & 类型 & 经济含义 & 状态 \\
\midrule
$i$ & 企业索引 & 不同企业 & 外生索引 \\
$r\in\{I,E,T,A\}$ & 路线索引 & 内部生产、委托生产、转让、放弃 & 外生索引 \\
$M\in\{0,1\}$ & 制度变量 & 是否法律上开放 MAH 外部 retained route & 外生制度 \\
$\eta\in[0,1]$ & 连续制度强度 & post-MAH implementation intensity & 外生制度强度 \\
$a_i$ & 企业特征 & 研发能力或研发生产率 & 外生异质性 \\
$k_i$ & 企业特征 & 内部生产能力 & 外生异质性 \\
$h_i^I\in\{0,1\}$ & 可行性指标 & 企业是否可行地走内部生产路线 & 外生异质性 \\
$q_i^E\in\{0,1\}$ & 可行性指标 & 企业是否具备进入合格委托生产安排的技术与监管资格；不表示已经完成匹配 & 外生异质性 \\
$s_i\in[0,1]$ & 项目基础成功率 & 项目自身从“值得规划”到“可能被实现”的基础成功部分 & 外生异质性 \\
$\chi_i^r$ & 路线实施因子 & 给定路线后，制度、执行和支持条件把项目真正落地的可能性 & 路线特定 \\
$\zeta_i^r=s_i\chi_i^r$ & 复合实现概率 & route-specific observed realization probability & 内生于路线环境 \\
\bottomrule
\end{tabularx}

\medskip
\noindent\begin{tabularx}{\textwidth}{p{0.13\textwidth} p{0.17\textwidth} X p{0.14\textwidth}}
\toprule
符号 & 类型 & 经济含义 & 状态 \\
\midrule
$R_i^r$ & 回报对象 & 路线 $r$ 成功实现时的事件性商业化回报 & 给定回报环境 \\
$v$ & continuation value & 新增一个 retained product 进入企业 stock 的边际延续价值 & 内生值函数对象 \\
$p$ & 有效支持成本 & 合格 CMO 支持的有效价格或支付成本 & 市场内生 \\
$\tau_i^E(M,\eta)$ & 路线摩擦 & MAH 可改善的 external route friction & 制度相关 \\
$\mu_i^E$ & 剩余 holder 负担 & 质量责任、协调、监督、技术转移等未消失成本 & 外生或缓变 \\
$S_i-\tau_i^T$ & 转让净回报 & 非 retained transfer/sale/out-licensing outside option & 外生基准 \\
$\kappa$ & 成本参数 & R\&D effort 的二次成本斜率 & 外生参数 \\
$H$ & 分布 & 企业异质性的横截面分布 & 外生分布 \\
$D^B(p)$ & 背景需求 & 模型外既有 CMO 支持需求 & 外生函数 \\
$S^{CMO}(p)$ & 支持供给 & 合格 CMO 支持供给函数 & 外生供给函数 \\
$p^*(M,\eta)$ & 均衡成本 & CMO 支持市场清算出的有效支持成本 & 内生结果 \\
\bottomrule
\end{tabularx}
\endgroup

\subsection{核心假设}

\begin{assumption}[部分均衡边界]
除合格 CMO 支持市场以外，其余背景环境都看作给定。模型不闭合产品市场、
研发投入市场、自由进入退出和行业稳态分布。
\end{assumption}

\begin{assumption}[路线可行性]
内部路线仅在 $h_i^I=1$ 时可行；外部 retained route 仅在
$M=1$ 且 $q_i^E=1$ 时可行。
\end{assumption}

\begin{assumption}[复合实现概率分解]
对任何产生观测结果的路线 $r$，
\[
\zeta_i^r(M,\eta,p)=s_i\chi_i^r(M,\eta,p), \qquad 0\le \zeta_i^r\le 1.
\]
其中 $s_i$ 是项目基础成功部分，$\chi_i^r$ 是 route-specific implementation
部分。
\end{assumption}

\begin{assumption}[实施强度的方向]
对 retained external route，$\partial \chi_i^E/\partial \eta\ge 0$，
$\partial \tau_i^E/\partial \eta\le 0$。也就是说，实施强度越高，
项目越容易通过外部路线落地，同时制度相关摩擦越低。
\end{assumption}

\begin{assumption}[R\&D 成本]
企业选择研发强度 $x_i\ge 0$，成本函数为
\[
C_i^R(x_i)=\frac{\kappa}{2}x_i^2,\qquad \kappa>0.
\]
这保证一阶条件简单、二阶条件严格凹。
\end{assumption}

\begin{assumption}[EV1 路线冲击]
路线选择阶段存在独立同分布的 Type-I extreme value 冲击，
尺度参数为 $\sigma>0$。
\end{assumption}

\begin{assumption}[CMO 支持市场的单调性]
$D^B(p)$ 连续且弱下降，$S^{CMO}(p)$ 连续且严格上升。
对每个企业，$G_i^E$ 对 $p$ 弱下降，因此由此诱导的
$x_i^*(M,\eta,p)$ 和 $P_i(E\mid M,\eta,p)$ 也弱下降。
\end{assumption}

\paragraph{为什么这些假设足够。}
这些假设并不要求我们解释所有宏观均衡，只要求我们把与 MAH
机制最直接相关的一块市场闭合起来。这样模型既有明确经济含义，
又不会做出超出数据与论文目标的夸张识别声明。

\paragraph{本节小结。}
这个模型最关键的对象有四组：企业异质性 $(a_i,k_i)$、制度状态 $(M,\eta)$、
路线实现对象 $\zeta_i^r=s_i\chi_i^r$、以及 CMO 支持成本 $p^*$。

\section{逐期 timing、路线集合与四类 payoff}

\subsection{逐期 timing}

把企业决策拆成一个足够清晰、又不过分复杂的顺序：

\begin{enumerate}[leftmargin=2em]
    \item 期初，企业带着既有产品存量 $n_i$、研发能力 $a_i$、
    内部生产能力 $k_i$ 等状态进入本期。
    \item 企业在看到当前制度环境 $(M,\eta)$ 后，选择原研研发强度 $x_i$。
    \item 研发强度决定“到达路线规划阶段的项目到达率”：
    \[
    \lambda_i^{lat}=a_i x_i.
    \]
    这里“lat”不是实验室最早想法，而是已经走到值得进入商业化路线规划阶段的
    有效项目流量。
    \item 如果项目到达了路线规划阶段，企业比较可行路线集合中的各个 payoff。
    \item 给定全体企业的外部路线使用倾向，CMO 支持市场清算出有效支持成本
    $p^*(M,\eta)$。
    \item 实际被观测到的 realized output 还要乘上 route-specific
    realization probability $\zeta_i^r$。
    \item retained route 成功实现后，企业下期 stock 增加；非 retained
    route 不增加 stock。
\end{enumerate}

\subsection{路线集合}

条件于项目已经到达路线规划阶段，企业的可行路线集合是
\begin{equation}
\Rset_i(M)=
\{A,T\}\cup \{I:h_i^I=1\}\cup \{E:M=1,\ q_i^E=1\}.
\label{eq:route_set}
\end{equation}

四条路线分别是：

\begin{enumerate}[leftmargin=2em]
    \item $I$：企业自己保留 holder 身份并内部生产；
    \item $E$：企业保留 holder 身份，但委托外部合格生产主体生产；
    \item $T$：转让、出售或 non-retained out-licensing；
    \item $A$：放弃或延迟。
\end{enumerate}

\subsection{四类 payoff}

\paragraph{第一步：先把实现概率与 continuation value 讲清楚。}
保留路线 $I,E$ 的共同特征是：只有项目真正被实现并形成 retained outcome，
企业才会拿到当期商业化事件回报，并把产品或证书加到自己的 stock 里。
因此，保留路线的 payoff 自然写成
\[
\zeta_i^r(R_i^r+v)-\text{cost},
\qquad r\in\{I,E\}.
\]

其中：

\begin{enumerate}[leftmargin=2em]
    \item $R_i^r$ 是事件性商业化回报；
    \item $v$ 是多保留一个产品进入 stock 的边际 continuation value；
    \item $\zeta_i^r=s_i\chi_i^r$ 是这条路线真正被落地并被观测到的复合实现概率。
\end{enumerate}

\paragraph{第二步：分别写每条路线。}

\begin{align}
G_i^I(M,\eta,p)
&=
\zeta_i^I(M,\eta,p)\bigl(R_i^I+v\bigr)-C^I(k_i),
\label{eq:GI}\\
G_i^E(M,\eta,p)
&=
\zeta_i^E(M,\eta,p)\bigl(R_i^E+v\bigr)
-p-\tau_i^E(M,\eta)-\mu_i^E,
\label{eq:GE}\\
G_i^T(M,\eta,p)
&=
S_i-\tau_i^T,
\label{eq:GT}\\
G_i^A(M,\eta,p)
&=0.
\label{eq:GA}
\end{align}

\paragraph{为什么 $T$ 不写成加 $v$ 的形式。}
$T$ 是 non-retained route。企业可以得到一笔净转让收益
$S_i-\tau_i^T$，但不会把产品保留在自己的 stock 里，因此没有 $+v$。
这一步非常重要，因为很多读者会把“委托生产”与“转让”混为一谈。
在这个模型里两者完全不同：

\begin{enumerate}[leftmargin=2em]
    \item $E$：holder 仍然是原企业，产品留在企业资产里；
    \item $T$：holder 不再是原企业，产品不进入企业 future stock。
\end{enumerate}

\paragraph{内部生产成本的一个便于直觉的例子。}
为了理解 $k_i$ 的作用，可以采用
\begin{equation}
C^I(k_i)=F_I+\frac{\bar c_I}{k_i},\qquad k_i>0.
\label{eq:CI_example}
\end{equation}
这里 $F_I$ 是内部生产的固定组织负担，$\bar c_I/k_i$ 表示内部生产能力越强，
单位项目的内部落地越便宜。重要的是方向：
\[
\frac{\partial C^I(k_i)}{\partial k_i}<0.
\]

\subsection{离散的 route-set expansion}

现在可以非常直接地看出 MAH 的离散效应。

\begin{align*}
\Rset_i(0)&=\{A,T\}\cup\{I:h_i^I=1\},\\
\Rset_i(1)&=\{A,T\}\cup\{I:h_i^I=1\}\cup\{E:q_i^E=1\}.
\end{align*}

因此，$M$ 的核心作用首先不是“让所有现有路线都更赚钱”，
而是\textbf{给部分企业多增加了一条原来根本不存在的 retained route}。
这就是 route-set expansion。

\paragraph{本节小结。}
本节最重要的式子是四条路线的 payoff。保留路线拿到的是
$\zeta_i^r(R_i^r+v)-\text{cost}$，非保留转让路线没有 $+v$。
MAH 的第一个作用是把 $E$ 路线从“不可选”变成“可选”。

\section{EV1 logit 与 inclusive value：一步一步推导}

\subsection{为什么需要随机效用}

如果企业永远机械地选择 payoff 最大的路线，那么模型会得到一个硬阈值：
谁的 $G_i^E$ 稍微高一点，谁就全部走 $E$；稍微低一点，就完全不走。
这在理论上可以，但在数据里往往太硬。现实中，即使面对同样的可观察 payoff，
不同企业或项目也会因为细微的管理、匹配、偏好差异做出不同路线选择。

因此我们写：
\begin{equation}
U_i^r = G_i^r + \varepsilon_i^r,
\label{eq:latent_utility}
\end{equation}
其中 $\varepsilon_i^r$ 是独立同分布的 EV1 冲击，分布函数为
\begin{equation}
F(\varepsilon)=\exp\left[-\exp\left(-\frac{\varepsilon}{\sigma}\right)\right],
\qquad \sigma>0.
\label{eq:ev1_cdf}
\end{equation}

对应密度函数是
\begin{equation}
f(\varepsilon)
=
\frac{1}{\sigma}
\exp\left(-\frac{\varepsilon}{\sigma}\right)
\exp\left[-\exp\left(-\frac{\varepsilon}{\sigma}\right)\right].
\label{eq:ev1_pdf}
\end{equation}

\subsection{路线选择概率的推导}

企业选择路线 $r$ 的概率是
\[
P_i(r)=\Prb\bigl(U_i^r\ge U_i^\ell,\ \forall \ell\in\Rset_i(M)\bigr).
\]

把 $U_i^r=G_i^r+\varepsilon_i^r$ 代进去，可得
\[
P_i(r)
=
\Prb\bigl(
\varepsilon_i^\ell
\le
\varepsilon_i^r + G_i^r-G_i^\ell,\ \forall \ell\neq r
\bigr).
\]

对 $\varepsilon_i^r$ 积分，有
\begin{align}
P_i(r)
&=
\int_{-\infty}^{\infty}
\prod_{\ell\neq r}
F\!\left(\varepsilon + G_i^r-G_i^\ell\right)
f(\varepsilon)\dd \varepsilon.
\label{eq:logit_step1}
\end{align}

把 EV1 的 $F(\cdot)$ 和 $f(\cdot)$ 代入后，乘积项可以整理成
\begin{align}
P_i(r)
&=
\int_{-\infty}^{\infty}
\frac{1}{\sigma}
\exp\left(-\frac{\varepsilon}{\sigma}\right)
\exp\left[
-\exp\left(-\frac{\varepsilon}{\sigma}\right)
\sum_{\ell\in\Rset_i(M)}
\exp\left(\frac{G_i^\ell-G_i^r}{\sigma}\right)
\right]
\dd \varepsilon.
\label{eq:logit_step2}
\end{align}

现在做变量替换：
\[
y=\exp\left(-\frac{\varepsilon}{\sigma}\right),\qquad
\dd y = -\frac{1}{\sigma}\exp\left(-\frac{\varepsilon}{\sigma}\right)\dd\varepsilon.
\]
于是积分变成
\begin{align}
P_i(r)
&=
\int_{0}^{\infty}
\exp\left[
-y
\sum_{\ell\in\Rset_i(M)}
\exp\left(\frac{G_i^\ell-G_i^r}{\sigma}\right)
\right]
\dd y \notag\\
&=
\left[
\sum_{\ell\in\Rset_i(M)}
\exp\left(\frac{G_i^\ell-G_i^r}{\sigma}\right)
\right]^{-1}.
\label{eq:logit_step3}
\end{align}

再把 $G_i^r$ 移回来，就得到标准 logit 公式：
\begin{equation}
P_i(r\mid M,\eta,p)
=
\frac{\exp(G_i^r/\sigma)}
{\sum_{\ell\in\Rset_i(M)}\exp(G_i^\ell/\sigma)}.
\label{eq:logit_prob}
\end{equation}

\subsection{inclusive value 的推导}

定义一项有效项目进入路线选择阶段时的期望最优价值：
\[
\E\left[\max_{r\in\Rset_i(M)}U_i^r\right].
\]

在 EV1 设定下有标准结果：
\[
\E\left[\max_{r}U_i^r\right]
=
\sigma\log\sum_{r}\exp(G_i^r/\sigma)+C_E\sigma,
\]
其中 $C_E$ 是 Euler 常数。由于 $C_E\sigma$ 对所有企业和路线都只是一个
加法常数，我们在模型里把它吸收到归一化里，定义
\begin{equation}
\Gamma_i(M,\eta,p)
=
\sigma\log\sum_{r\in\Rset_i(M)}\exp(G_i^r/\sigma).
\label{eq:inclusive_value}
\end{equation}

\paragraph{为什么 $\Gamma_i$ 很重要。}
$\Gamma_i$ 不是“某条路线”的价值，而是“一旦项目到达路线规划阶段，
企业平均而言能拿到的最优商业化机会价值”。它把所有路线选择可能性压缩成
一个标量，因此非常适合作为研发阶段的前瞻回报对象。

\subsection{一个关键导数}

对任意路线 payoff $G_i^r$，对式 \eqref{eq:inclusive_value} 求导：
\begin{align}
\frac{\partial \Gamma_i}{\partial G_i^r}
&=
\sigma
\frac{1}{\sum_{\ell}\exp(G_i^\ell/\sigma)}
\cdot
\frac{1}{\sigma}\exp(G_i^r/\sigma)\notag\\
&=
\frac{\exp(G_i^r/\sigma)}
{\sum_{\ell}\exp(G_i^\ell/\sigma)}
=
P_i(r\mid M,\eta,p).
\label{eq:Gamma_derivative}
\end{align}

这个式子有极强的直觉含义：

\begin{quote}
如果某条路线在最优选择里本来就几乎不会被采用，那么改善它的 payoff
对 $\Gamma_i$ 的提升就很小；如果某条路线本来就很 relevant，
改善它的 payoff 对 $\Gamma_i$ 的提升就大。
\end{quote}

\subsection{确定性极限}

当 $\sigma\to 0$ 时，
\[
\Gamma_i(M,\eta,p)\to \max_{r\in\Rset_i(M)} G_i^r(M,\eta,p).
\]
这说明 deterministic max 不是另一个完全不同的模型，而是 logit 模型的极限。
但在本讲义里，基准形式就是 logit，因为后面 route share、observed share、
log odds 的解释都依赖这个平滑形式。

\paragraph{本节小结。}
EV1 假设给出两个核心结果：一是 logit 选择概率，二是 inclusive value
$\Gamma_i$。更重要的是，$\partial\Gamma_i/\partial G_i^r=P_i(r)$，
这让后面的比较静态非常透明。

\section{R\&D 选择：FOC、SOC 与最优研发强度}

\subsection{有效项目到达率}

按照用户给定的统一模型，\textbf{$a_i x_i$ 是进入路线规划阶段的项目到达率}：
\begin{equation}
\lambda_i^{lat}=a_i x_i.
\label{eq:latent_arrival}
\end{equation}

这个写法的含义是：

\begin{enumerate}[leftmargin=2em]
    \item $a_i$ 衡量企业把研发投入转成有效项目机会的能力；
    \item $x_i$ 是企业自己选择的研发强度；
    \item 两者相乘，得到进入“后续商业化路线选择”的项目流量。
\end{enumerate}

\subsection{研发问题}

在给定 $(M,\eta,p)$ 的情况下，企业的研发问题写成
\begin{equation}
\max_{x_i\ge 0}
\left\{
\beta a_i x_i \Gamma_i(M,\eta,p)
\;-\;
\frac{\kappa}{2}x_i^2
\right\}.
\label{eq:rd_problem}
\end{equation}

这里的逻辑非常直接：

\begin{enumerate}[leftmargin=2em]
    \item 左边第一项 $\beta a_i x_i \Gamma_i$ 是“有效项目到达率”
    乘上“每个有效项目的前瞻机会价值”；
    \item 第二项是研发强度的凸成本。
\end{enumerate}

\subsection{一阶条件 FOC}

对 \eqref{eq:rd_problem} 关于 $x_i$ 求导：
\begin{equation}
\frac{\partial}{\partial x_i}
\left[
\beta a_i x_i \Gamma_i
-\frac{\kappa}{2}x_i^2
\right]
=
\beta a_i\Gamma_i - \kappa x_i.
\label{eq:foc_step}
\end{equation}

令一阶条件等于零，得到
\begin{equation}
\beta a_i\Gamma_i(M,\eta,p)-\kappa x_i=0.
\label{eq:foc}
\end{equation}

所以最优研发强度是
\begin{equation}
x_i^*(M,\eta,p)
=
\frac{\beta a_i}{\kappa}\Gamma_i(M,\eta,p).
\label{eq:xstar_offeq}
\end{equation}

\subsection{二阶条件 SOC}

对 \eqref{eq:rd_problem} 再求一次导数：
\[
\frac{\partial^2}{\partial x_i^2}
\left[
\beta a_i x_i \Gamma_i
-\frac{\kappa}{2}x_i^2
\right]
=
-\kappa<0.
\]

因此目标函数严格凹，\eqref{eq:xstar_offeq} 给出的就是唯一最优解。

\subsection{把最优研发强度代回去}

最优有效项目到达率就是
\begin{equation}
\lambda_i^{lat,*}(M,\eta,p)
=
a_i x_i^*(M,\eta,p)
=
\frac{\beta a_i^2}{\kappa}\Gamma_i(M,\eta,p).
\label{eq:lambda_lat_star}
\end{equation}

这个式子非常重要，因为它说明：

\begin{quote}
原研项目流量不是单独由研发能力 $a_i$ 决定，也不是单独由商业化机会价值
$\Gamma_i$ 决定，而是二者共同决定。
\end{quote}

\subsection{观测到的 realized output}

潜在项目流量 $\lambda_i^{lat,*}$ 和最终观测到的 output 不是一回事。
路线选择和路线实现都要继续筛选。定义路线 $r$ 的观测实现输出为
\begin{equation}
Y_i^{obs,r}(M,\eta,p)
=
a_i x_i^*(M,\eta,p)
P_i(r\mid M,\eta,p)
\zeta_i^r(M,\eta,p).
\label{eq:obs_route_output}
\end{equation}

于是企业的总观测 realized output 是
\begin{equation}
Y_i^{obs}(M,\eta,p)
=
a_i x_i^*(M,\eta,p)
\sum_{r\in\Rset_i(M)}
P_i(r\mid M,\eta,p)\zeta_i^r(M,\eta,p).
\label{eq:obs_total_output}
\end{equation}

\paragraph{直觉。}
式 \eqref{eq:obs_total_output} 有三层乘法：

\begin{enumerate}[leftmargin=2em]
    \item 有多少有效项目到达路线规划阶段；
    \item 这些项目有多大概率走某条路线；
    \item 走了那条路线后，有多大概率真的实现并被我们观测到。
\end{enumerate}

所以数据里的 approvals、holder--producer split、observed original-drug
records，本质上都在看第三层对象，而不是直接看到最上游的“科研灵感”。

\paragraph{本节小结。}
最优研发强度是
$x_i^*=\frac{\beta a_i}{\kappa}\Gamma_i$，唯一性来自严格凹的二次成本。
观测 output 则必须进一步乘上路线选择概率和复合实现概率。

\section{CMO 支持市场：背景需求、存在唯一性与 IFT}

\subsection{为什么要单独闭合这一块市场}

如果所有企业都想走 retained external route，那么他们对合格 CMO
支持的需求就会上升。若支持资源稀缺，价格会上升，改革带来的好处会被部分抵消。
这正是 MAH 机制里最有经济内容、也最应该被保留的价格反馈。

\subsection{总需求函数}

按照用户指定的结构，CMO 总需求是
\begin{equation}
D^{CMO}(p;M,\eta)
=
D^B(p)
+\int a_i x_i^*(M,\eta,p)P_i(E\mid M,\eta,p)\dd H(i).
\label{eq:cmo_total_demand}
\end{equation}

这里特别注意两点。

\begin{enumerate}[leftmargin=2em]
    \item 第一项 $D^B(p)$ 是背景需求，表示模型外已有的 CMO 支持需求。
    \item 第二项是本模型企业在\textbf{路线规划阶段}对 external support 的需求。
    \textbf{这里不能再乘 $\zeta_i^E$。}因为这一步是“要不要为可能的
    external route 做支持安排”的需求，不是最终 realized output。
\end{enumerate}

\subsection{供给函数与市场清算}

设合格 CMO 支持供给为
\begin{equation}
S^{CMO}(p),
\qquad
\frac{\dd S^{CMO}(p)}{\dd p}>0.
\label{eq:cmo_supply}
\end{equation}

于是市场清算条件是
\begin{equation}
S^{CMO}(p)=D^{CMO}(p;M,\eta).
\label{eq:market_clearing}
\end{equation}

均衡有效支持成本记为 $p^*(M,\eta)$。

\subsection{存在与唯一性}

定义
\begin{equation}
F(p;M,\eta)
\equiv
S^{CMO}(p)-D^{CMO}(p;M,\eta).
\label{eq:F_function}
\end{equation}

\begin{proposition}[CMO 支持市场的存在与唯一性]
如果：
\begin{enumerate}[label=(\roman*),leftmargin=2em]
    \item $S^{CMO}(p)$ 连续且严格上升；
    \item $D^B(p)$ 连续且弱下降；
    \item 对几乎所有企业，$x_i^*(M,\eta,p)$ 与
    $P_i(E\mid M,\eta,p)$ 对 $p$ 连续且弱下降；
    \item 存在两个价格点 $\underline p<\bar p$，使得
    $F(\underline p;M,\eta)\le 0$ 且 $F(\bar p;M,\eta)\ge 0$，
\end{enumerate}
那么至少存在一个清算价格 $p^*(M,\eta)$。如果进一步
$F(p;M,\eta)$ 严格上升，那么这个清算价格唯一。
\end{proposition}

\paragraph{证明思路。}
因为 $F(p;M,\eta)$ 连续，且在 $\underline p$ 和 $\bar p$ 两端异号，
由介值定理，存在 $p^*$ 使得 $F(p^*;M,\eta)=0$。如果 $F$ 严格上升，
那么 $F$ 只能穿过零点一次，因此均衡唯一。

\paragraph{经济解释。}
唯一性本质上要求一件事：价格越高，支持供给越多，而企业想走 external
route 的需求越少。只要这两个方向足够稳定，市场就不会出现多个互不相干的
清算价格。

\subsection{IFT：实施强度上升时，\texorpdfstring{$p^*$}{均衡支持成本}如何变化}

现在看 post-MAH implementation intensity $\eta$ 的作用。
在均衡点上有
\[
F\bigl(p^*(M,\eta);M,\eta\bigr)=0.
\]

如果 $F$ 在均衡点可微且 $F_p\neq 0$，由隐函数定理，
\begin{equation}
\frac{\dd p^*(M,\eta)}{\dd \eta}
=
-\frac{F_\eta(p^*;M,\eta)}{F_p(p^*;M,\eta)}.
\label{eq:ift}
\end{equation}

由于 $F_p>0$，所以符号由 $F_\eta$ 决定。若更强的实施强度会提高
企业走 $E$ 路线的意愿，从而使
\[
\frac{\partial D^{CMO}(p;M,\eta)}{\partial \eta}>0,
\]
那么
\[
F_\eta(p;M,\eta)
=
-\frac{\partial D^{CMO}(p;M,\eta)}{\partial \eta}<0,
\]
于是
\begin{equation}
\frac{\dd p^*(M,\eta)}{\dd \eta}>0.
\label{eq:dpdeta_positive}
\end{equation}

\paragraph{直觉。}
这就是说：实施越顺畅，越多项目愿意规划 external route，
对合格 CMO 支持的需求越大，支持成本就越容易上升。
这就是价格反馈。

\paragraph{本节小结。}
CMO 市场只需要闭合一条市场清算式：
$S^{CMO}(p)=D^B(p)+\int a_i x_i^*P_i(E)\dd H$。
只要供给上升、需求下降、边界条件满足，均衡存在且唯一。IFT 则告诉我们，
$\eta$ 上升很可能推高 $p^*$。

\section{离散 route-set expansion 与 \texorpdfstring{$\eta$}{实施强度}比较静态}

\subsection{\texorpdfstring{$M$}{制度开放}的离散效应：新增一条路线}

在给定其他 payoff 不变时，pre-MAH 与 post-MAH 的 inclusive value 可写成
\begin{align}
\Gamma_i(0,0,p)
&=
\sigma\log
\sum_{r\in\Rset_i(0)}
\exp(G_i^r/\sigma),
\label{eq:gamma_pre}\\
\Gamma_i(1,\eta,p)
&=
\sigma\log
\left[
\sum_{r\in\Rset_i(0)}
\exp(G_i^r/\sigma)
+
\exp(G_i^E(1,\eta,p)/\sigma)
\right].
\label{eq:gamma_post}
\end{align}

因此，只要 $E$ 路线在 post-MAH 下可行且 payoff 不是负无穷，就有
\begin{equation}
\Gamma_i(1,\eta,p)\ge \Gamma_i(0,0,p).
\label{eq:route_set_expansion}
\end{equation}

\paragraph{解释。}
多一条可选路线，最优选择集合只会变大，不会变小。这是 route-set expansion
最直接的数学表达。

\subsection{\texorpdfstring{$\eta$}{实施强度}的直接效应}

保持 $p$ 固定时，只有 external route payoff 直接受 $\eta$ 影响：
\begin{equation}
\frac{\partial G_i^E}{\partial \eta}\Big|_{p}
=
(R_i^E+v)\frac{\partial \zeta_i^E}{\partial \eta}
-\frac{\partial \tau_i^E}{\partial \eta}.
\label{eq:dGE_deta_partial}
\end{equation}

因为 $\zeta_i^E=s_i\chi_i^E$，所以
\[
\frac{\partial \zeta_i^E}{\partial \eta}
=
s_i\frac{\partial \chi_i^E}{\partial \eta}.
\]
于是
\begin{equation}
\frac{\partial G_i^E}{\partial \eta}\Big|_{p}
=
s_i(R_i^E+v)\frac{\partial \chi_i^E}{\partial \eta}
-\frac{\partial \tau_i^E}{\partial \eta}.
\label{eq:dGE_deta_partial2}
\end{equation}

在常见方向下，
$\partial\chi_i^E/\partial\eta\ge 0$，
$\partial\tau_i^E/\partial\eta\le 0$，
所以右边通常为正。

\subsection{\texorpdfstring{$\eta$}{实施强度}对 inclusive value 的作用}

由前面的关键导数 \eqref{eq:Gamma_derivative}，
\begin{equation}
\frac{\partial \Gamma_i}{\partial \eta}\Big|_{p}
=
P_i(E\mid M,\eta,p)\cdot
\frac{\partial G_i^E}{\partial \eta}\Big|_{p}.
\label{eq:dGamma_deta_partial}
\end{equation}

这说明：

\begin{quote}
实施强度上升并不会平均地提高所有企业的机会价值。只有那些本来就会把
$E$ 路线当回事的企业，$\Gamma_i$ 才会有较大改善。
\end{quote}

\subsection{\texorpdfstring{$\eta$}{实施强度}对研发强度的作用}

由
$x_i^*=\frac{\beta a_i}{\kappa}\Gamma_i$，
可得
\begin{equation}
\frac{\partial x_i^*}{\partial \eta}\Big|_{p}
=
\frac{\beta a_i}{\kappa}
\frac{\partial \Gamma_i}{\partial \eta}\Big|_{p}
=
\frac{\beta a_i}{\kappa}
P_i(E\mid M,\eta,p)
\frac{\partial G_i^E}{\partial \eta}\Big|_{p}.
\label{eq:dx_deta_partial}
\end{equation}

\subsection{总效应：把价格反馈也算进去}

在均衡里，$p=p^*(M,\eta)$ 也会跟着 $\eta$ 动。于是
\begin{equation}
\frac{\dd G_i^E}{\dd \eta}
=
\frac{\partial G_i^E}{\partial \eta}
+
\frac{\partial G_i^E}{\partial p}
\frac{\dd p^*}{\dd \eta}.
\label{eq:total_GE_eta}
\end{equation}

而
\begin{equation}
\frac{\partial G_i^E}{\partial p}
=
(R_i^E+v)\frac{\partial \zeta_i^E}{\partial p} -1.
\label{eq:dGE_dp}
\end{equation}

通常 $\partial \zeta_i^E/\partial p\le 0$，因为支持越贵、越稀缺，
路线实现越困难，因此
\[
\frac{\partial G_i^E}{\partial p}<0.
\]

于是总效应变成
\begin{equation}
\frac{\dd \Gamma_i}{\dd \eta}
=
P_i(E\mid M,\eta,p^*)
\left[
\frac{\partial G_i^E}{\partial \eta}
+
\frac{\partial G_i^E}{\partial p}
\frac{\dd p^*}{\dd \eta}
\right].
\label{eq:total_gamma_eta}
\end{equation}

进而
\begin{equation}
\frac{\dd x_i^*}{\dd \eta}
=
\frac{\beta a_i}{\kappa}
\frac{\dd \Gamma_i}{\dd \eta}.
\label{eq:total_x_eta}
\end{equation}

\subsection{attenuation 的含义}

若 $\dd p^*/\dd\eta>0$，而 $\partial G_i^E/\partial p<0$，
那么
\[
\frac{\partial G_i^E}{\partial p}\frac{\dd p^*}{\dd \eta}<0.
\]
这就是 attenuation：

\begin{quote}
更强的 MAH 实施强度一方面直接改善 external route，
另一方面又通过支持市场推高 $p^*$，从而吃掉一部分直接收益。
\end{quote}

\paragraph{什么时候可能看不到很强的总效应。}
哪怕 $\eta$ 提高了制度执行质量，如果某个企业：

\begin{enumerate}[leftmargin=2em]
    \item 本来就很少会选 $E$；
    \item 或者 $p^*$ 上升得很快；
    \item 或者 $\mu_i^E$ 很大；
    \item 或者转让 outside option 很强，
\end{enumerate}
那么最终 $\dd x_i^*/\dd\eta$ 也可能不大。

\paragraph{本节小结。}
$M$ 主要带来离散的 route-set expansion，$\eta$ 主要带来连续的实施改善。
但一旦把 $p^*$ 的价格反馈算进去，$\eta$ 的好处会被 attenuate。

\section{observed output、observed share 与 log odds 的使用条件}

\subsection{路线层面的观测输出}

前面已经定义了
\[
Y_i^{obs,r}=a_i x_i^* P_i(r)\zeta_i^r.
\]
其中 external retained route 的观测输出是
\begin{equation}
Y_i^{obs,E}(M,\eta)
=
a_i x_i^*(M,\eta,p^*)
P_i(E\mid M,\eta,p^*)
\zeta_i^E(M,\eta,p^*).
\label{eq:obsE}
\end{equation}

\subsection{正确的 aggregate observed holder--producer share}

这是本讲义中特别容易被写错的一步。若我们研究的是
\textbf{观测到的 retained original-drug records 中 holder 与 producer 分离的份额}，
那么正确的 aggregate observed share 不是简单平均的
$P_i(E)$，也不是简单平均的 $P_i(E)\zeta_i^E$，
而必须按“有多少项目流量真正走到被观测的 realized output”
来加权。

令 $\Rset_i^{ret}(M)=\Rset_i(M)\cap\{I,E\}$。先分别定义 observed entrusted
flow 与全部 observed retained flow：
\begin{align}
N_E^{obs}(M,\eta)
&=
\int a_i x_i^*(M,\eta,p^*)P_i(E\mid M,\eta,p^*)
\zeta_i^E(M,\eta,p^*)\dd H(i),
\label{eq:observed_e_flow}\\
N_{ret}^{obs}(M,\eta)
&=
\int a_i x_i^*(M,\eta,p^*)
\sum_{r\in\Rset_i^{ret}(M)}
P_i(r\mid M,\eta,p^*)\zeta_i^r(M,\eta,p^*)\dd H(i).
\label{eq:observed_retained_flow}
\end{align}
于是正确的 aggregate share 是
\begin{equation}
Share^{HP}(M,\eta)
=
\frac{N_E^{obs}(M,\eta)}{N_{ret}^{obs}(M,\eta)}.
\label{eq:aggregate_share}
\end{equation}

\paragraph{为什么这个权重要用 $a_i x_i^*P_i\zeta_i$。}
因为真正进入数据的人次，是由下面四件事共同决定的：

\begin{enumerate}[leftmargin=2em]
    \item 企业研发出了多少有效项目；
    \item 这些项目有多少走路线 $r$；
    \item 走了这条路线后有多少真正被实现；
    \item 这些实现结果是否进入我们定义的样本口径。
\end{enumerate}

因此，观测 share 的本质是 realized-flow share，而不是 choice-probability
share。

\subsection{两路线比较下的 log odds}

如果我们在一个\textbf{可比的 retained 子样本}里只比较 $I$ 与 $E$，
那么有
\begin{equation}
s_i^E
=
\frac{\exp(G_i^E/\sigma)}
{\exp(G_i^I/\sigma)+\exp(G_i^E/\sigma)}.
\label{eq:two_route_share}
\end{equation}

于是
\begin{equation}
\log\frac{s_i^E}{1-s_i^E}
=
\frac{G_i^E-G_i^I}{\sigma}.
\label{eq:logodds_basic}
\end{equation}

再把 payoff 代进去：
\begin{align}
G_i^E-G_i^I
&=
\zeta_i^E(R_i^E+v)-p^*-\tau_i^E-\mu_i^E
\notag\\
&\quad
-\Bigl[\zeta_i^I(R_i^I+v)-C^I(k_i)\Bigr].
\label{eq:payoff_gap}
\end{align}

所以 log odds 反映的是一整个\textbf{相对有效路线楔子}，
不是某一个 primitive parameter。

\subsection{什么时候 log odds 可以被认真解释}

要把 log odds 用成一个有意义的经验对象，至少要满足以下条件。

\begin{enumerate}[leftmargin=2em]
    \item 比较的确实是可比 retained route。也就是分子分母都对应
    retained internal 与 retained external，而不是把 transfer、abandonment
    混进去。
    \item 样本口径与理论口径一致。若数据定义的是 observed retained
    original-drug records，就必须用对应的 realized-output 权重。
    \item 除 external-route wedge 以外，其他系统性差异要么稳定，要么被控制。
    否则 log odds 变化可能混入 $C^I(k_i)$、$R_i^I$、$R_i^E$、$\mu_i^E$、
    样本构成变化等因素。
    \item pre 和 post 必须存在法律和经济上可比的 external route，
    否则 pre--post log odds 本身就不合法。
\end{enumerate}

\subsection{为什么 pre-MAH 零份额不能机械代入 log odds}

如果 pre-MAH 时 $E$ 路线根本不存在，那么
\[
\Rset_i(0)=\{A,T\}\cup\{I:h_i^I=1\},
\]
此时不存在一个真正可比的 pre-reform entrusted retained share。
机械地把
\[
s_{pre}^E=0
\]
代入
\[
\log\frac{s_{pre}^E}{1-s_{pre}^E}
\]
会得到 $\log(0)$，这是未定义的。

因此，在 MAH 这类 route-set expansion 问题里，更稳妥的做法通常是：

\begin{enumerate}[leftmargin=2em]
    \item 把 pre-MAH 解释成“没有 $E$ 路线”的边界状态；
    \item 在 post-MAH 内部利用 implementation intensity $\eta$ 的横截面差异；
    \item 若一定要构造 pre--post log odds，则只能把 pseudo count 当成
    敏感性分析，而不能当基准事实。
\end{enumerate}

\subsection{post-MAH implementation intensity 的 log odds 规格}

在更适合本模型的 post-MAH 横截面里，可以写
\begin{equation}
\log\frac{s_\ell^{HP}}{1-s_\ell^{HP}}
=
\alpha_0+\alpha_1\eta_\ell+\bm{Z}_\ell'\bm{\alpha}+u_\ell,
\label{eq:logodds_eta}
\end{equation}
其中 $\ell$ 表示地区、年份、剂型、治疗领域或企业类型等 implementation cell。

这个回归的合理解释是：

\begin{quote}
$\alpha_1$ 反映的是 implementation intensity 与 relative external-route
attractiveness 之间的关系。
\end{quote}

它\textbf{不是}对 $\tau_i^E$、$\mu_i^E$、$\chi_i^E$ 中任何单一 primitive
参数的直接识别。

\paragraph{本节小结。}
观测 share 一定要按 $a_i x_i^*P_i\zeta_i$ 加权。log odds 只有在 retained route
可比、样本口径一致且零份额问题被妥善处理时才有明确含义。

\section{\texorpdfstring{$a_i$ 与 $k_i$}{研发能力与生产能力}异质性：哪些企业最会响应 MAH}

\subsection{\texorpdfstring{$a_i$}{研发能力}：研发能力高的企业为什么更敏感}

由
\[
x_i^*=\frac{\beta a_i}{\kappa}\Gamma_i
\]
直接可得
\begin{equation}
\frac{\partial x_i^*}{\partial a_i}
=
\frac{\beta}{\kappa}\Gamma_i
>0.
\label{eq:dx_da}
\end{equation}

所以研发能力更强的企业，对任意给定的商业化机会价值都更能把机会转成研发强度。
这意味着当 MAH 提高 $\Gamma_i$ 时，高 $a_i$ 企业的研发响应会更大。

\subsection{\texorpdfstring{$k_i$}{内部生产能力}：内部生产能力弱的企业为什么更依赖 external route}

看 retained external 与 retained internal 的相对吸引力：
\begin{equation}
\Delta_i^{EI}
\equiv
G_i^E-G_i^I.
\label{eq:delta_EI}
\end{equation}

若 $R_i^r$ 与 $\zeta_i^r$ 对 $k_i$ 的直接依赖不强，则
\begin{equation}
\frac{\partial \Delta_i^{EI}}{\partial k_i}
=
\frac{\partial C^I(k_i)}{\partial k_i}
<0.
\label{eq:delta_k}
\end{equation}

这个式子的意思是：

\begin{quote}
内部生产能力越强，内部路线越有吸引力，external route 相对吸引力越低。
\end{quote}

因此，MAH 机制最容易作用在“研发能力高、但内部生产能力弱”的企业上。

\subsection{转让 outside option 的作用}

别忘了路线 $T$。如果某企业的
\[
S_i-\tau_i^T
\]
很高，那么即便 external retained route 变得更顺畅，企业也可能更愿意转让，
而不是自己保留 holder 身份去商业化。这说明：

\begin{enumerate}[leftmargin=2em]
    \item external route 的 relevant set 不是所有企业；
    \item 高 $a_i$、低 $k_i$ 只是必要方向，不是充分条件；
    \item 还要看 $q_i^E$、$\mu_i^E$、$S_i-\tau_i^T$ 等对象。
\end{enumerate}

\subsection{一组清晰的横截面预测}

综合起来，模型预测 MAH 的 route-friction channel 对以下企业最强：

\begin{enumerate}[leftmargin=2em]
    \item $a_i$ 高：研发能力强，商业化机会价值更容易转化成研发投入；
    \item $k_i$ 低：内部生产路线昂贵，external retained route 更 relevant；
    \item $q_i^E=1$：能接触到合格 CMO 支持；
    \item $\mu_i^E$ 不太高：holder-side 剩余负担没有完全吃掉制度红利；
    \item $S_i-\tau_i^T$ 不太高：转让 outside option 不至于压过 retained route。
\end{enumerate}

\paragraph{本节小结。}
如果数据支持“高研发能力、低内部生产能力、且能接触合格 CMO 的企业，
holder--producer split 和原研结果响应更强”，这与 MAH 路线机制高度一致。

\section{Bellman 形式、expected stock transition 与 \texorpdfstring{$v$}{延续价值}}

\subsection{为什么还需要 Bellman 视角}

前面看起来已经够用了：有路线收益，有 logit，有研发一阶条件。
但如果要严谨说明为什么 retained route 会多一个 $+v$，就需要把 stock
的动态写出来。因为 $v$ 不是“凭空加上的奖励”，而是新增一个 retained
产品进入未来 stock 的边际价值。

\subsection{expected stock transition}

记企业已有 retained product stock 为 $n_i$，折旧率或退出率为 $\delta\in(0,1)$。
定义 retained indicator：
\[
\rho^I=\rho^E=1,\qquad
\rho^T=\rho^A=0.
\]

那么给定研发强度 $x_i$，下一期的\textbf{期望} stock 转移是
\begin{equation}
\E[n_i' \mid n_i,x_i,M,\eta]
=
(1-\delta)n_i
+
a_i x_i
\sum_{r\in\Rset_i(M)}
P_i(r\mid M,\eta,p^*)
\rho^r
\zeta_i^r(M,\eta,p^*).
\label{eq:stock_transition}
\end{equation}

这一步要慢慢看：

\begin{enumerate}[leftmargin=2em]
    \item 旧 stock 有一部分按 $1-\delta$ 留下来；
    \item 新增 stock 来自本期有效项目流量 $a_i x_i$；
    \item 但并不是所有项目都进 stock，只有 retained route 才会进；
    \item 即使选了 retained route，也还要乘上实现概率 $\zeta_i^r$。
\end{enumerate}

\subsection{Bellman 方程}

Bellman accounting 必须避免把新增 stock 的价值算两次。最清楚的做法是先写原始动态问题：
企业支付研发成本，下一期每个 route-planning project 产生一次路线事件回报，
并通过式 \eqref{eq:stock_transition} 改变下一期 stock。把仿射值函数代入并对路线
选择与实现取期望以后，得到下面的\textbf{约化 Bellman 方程}：
\begin{equation}
V_i(n_i;M,\eta)
=
\pi n_i
+\beta A_i(M,\eta)
+\beta(1-\delta)v n_i
+\max_{x_i\ge 0}\left\{
\beta a_i x_i\Gamma_i(M,\eta,p^*)
-\frac{\kappa}{2}x_i^2
\right\}.
\label{eq:bellman}
\end{equation}

这里 $\pi n_i$ 是已有 retained products 的当期流收益，
$\beta A_i+\beta(1-\delta)vn_i$ 是旧 stock 的下一期 continuation value，
而 $\Gamma_i$ 已经包括新增项目在各路线下的事件回报和 retained 路线新增的 $v$。

\paragraph{为什么这里不会 double count。}
关键规则是：

\begin{enumerate}[leftmargin=2em]
    \item $R_i^r$ 是项目本次成功商业化时的\textbf{事件性回报}；
    \item $v$ 是新增 retained asset 带来的\textbf{未来 stock 延续价值}；
    \item $\pi n_i$ 是\textbf{已有 stock} 的当期流收益。
    \item 一旦 $\Gamma_i$ 写成 $\zeta_i^r(R_i^r+v)-c_i^r$ 的 expected maximum，
    约化 Bellman 就不能再另外加一遍包含新增 stock 的
    $\beta\E[V_i(n_i')]$；否则同一个新增 $v$ 会被计算两次。
\end{enumerate}

因此式 \eqref{eq:bellman}只保留旧 stock 的 continuation term，并把新增项目的
全部边际价值放在 $\Gamma_i$ 中。这三类回报不是同一个东西，不能混为一谈。

\subsection{仿射值函数}

若环境平稳，写成
\begin{equation}
V_i(n_i;M,\eta)=A_i(M,\eta)+v\,n_i.
\label{eq:affine_value}
\end{equation}

则 $v$ 满足
\begin{equation}
v=\pi+\beta(1-\delta)v.
\label{eq:v_recursion}
\end{equation}

解得
\begin{equation}
v=\frac{\pi}{1-\beta(1-\delta)}.
\label{eq:v_closed_form}
\end{equation}

\paragraph{这个式子的经济意义。}
多一个 retained product 的价值，等于它本期带来的流收益 $\pi$，
再加上它在未来继续留在 stock 中的折现价值。只要
$\beta(1-\delta)<1$，这个值就是有限的。

\subsection{为什么 retained route payoff 要写成 \texorpdfstring{$(R_i^r+v)$}{事件回报加延续价值}}

有了 \eqref{eq:v_closed_form} 以后，保留路线 payoff 里的 $+v$
就完全有来历了：

\begin{quote}
内部生产或 retained external route 成功后，企业不只是拿到一次性事件回报，
还把产品保留下来，进入未来 stock，因此多拿一个 continuation value。
\end{quote}

\paragraph{本节小结。}
$v$ 来自 stock 的动态价值，而不是人为加上去的“奖励项”。
Bellman 写清楚以后，保留路线 payoff 中的 $(R_i^r+v)$ 就有了坚实基础。

\section{校准与识别：当前数据能识别什么，不能识别什么}

\subsection{为什么这里必须克制}

这个模型的经验任务不是“估计所有 primitive structural parameters”，
而是用当前可行的数据，去约束一组\textbf{复合对象}。原因很简单：
观测到的 approvals、holder--producer split、realized original-drug counts
都已经是下游结果，它们把研发能力、路线选择、实现概率、支持成本等东西
混在一起了。

\subsection{当前数据最自然约束的对象}

\begin{center}
\small
\begin{tabularx}{\textwidth}{>{\raggedright\arraybackslash}p{0.24\textwidth} >{\raggedright\arraybackslash}p{0.26\textwidth} X}
\toprule
数据矩 & 对应模型对象 & 识别边界 \\
\midrule
realized original-drug counts &
$\frac{\beta a_i^2}{\kappa}\Gamma_i$ 与 realization mapping 的组合 &
不能单独识别 $a_i$、$\kappa$、$\Gamma_i$、$\zeta_i^r$ \\
observed holder--producer split &
式 \eqref{eq:aggregate_share} 的 aggregate realized-flow 比率 &
不能单独识别 $\tau_i^E$、$\mu_i^E$、$\chi_i^E$、$p^*$ \\
不同 implementation cell 的 post-MAH log odds &
effective external-route wedge 相对 $\sigma$ 的变化 &
解释依赖可比 retained 子样本与控制变量 \\
CMO 容量或合格生产支持 proxy &
$p^*$ 的稀缺性与 attenuation 强弱 &
若没有好 proxy，就难以精确分离价格反馈 \\
\bottomrule
\end{tabularx}
\end{center}

\subsection{一个最简单的归一化思路}

假设我们只能看到 realized output 与 observed share，那么至少需要做一些归一化。
常见的做法包括：

\begin{enumerate}[leftmargin=2em]
    \item 归一化 $\kappa=1$；
    \item 归一化某类基准企业在 pre-MAH 下的 $\Gamma_i=1$；
    \item 把 $\sigma$ 看成离散选择平滑参数，先固定，再看复合楔子变化；
    \item 把 $s_i$ 与一部分 $\chi_i^r$ 合并到复合实现项里，而不是硬拆开。
\end{enumerate}

\subsection{建议的校准步骤}

\begin{enumerate}[leftmargin=2em]
    \item 先确定样本口径。比如只看 observed retained original-drug records，
    还是看 broader original-drug realized outcomes。
    \item 用 post-MAH 的 holder--producer split 或 implementation gradient
    去约束 effective external-route wedge 的相对变化。
    \item 用 realized original-drug counts 去约束
    $\frac{\beta a_i^2}{\kappa}\Gamma_i$ 这一类 composite scale。
    \item 若有 CMO 稀缺度、合格生产者数量、产能 proxy，再去评估 attenuation。
    \item 把 entry block 明确列为 future module。没有可信 firm-year
    research-oriented entry panel 时，不要把 entry 当成基准识别来源。
\end{enumerate}

\subsection{为什么 observed share 一定要用 realized-flow 权重}

如果把 observed share 错写成简单平均：
\[
\frac{\int P_i(E)\dd H(i)}{\int 1\dd H(i)},
\]
就等于把“一个几乎没有项目流量、几乎不会被观测到的企业”
和“一个拥有大量有效项目、而且实现率很高的企业”赋予同等权重。
这显然不对应数据里的 record share。正确做法是式
\eqref{eq:aggregate_share} 那样，用
$a_i x_i^*P_i\zeta_i$ 去加权。

\subsection{目前不能做出的强识别声明}

只靠当前这类 approval-side 或 realized-product-side 数据，不能认真声称：

\begin{enumerate}[leftmargin=2em]
    \item 单独识别法律合规成本；
    \item 单独识别搜索匹配成本；
    \item 单独识别 holder-side 监督负担；
    \item 单独识别项目真实成功概率；
    \item 单独识别合格 CMO 支持价格；
    \item 单独识别研究导向 entry 的真正成本分布。
\end{enumerate}

这些对象都以复合方式进入 payoff gap、route share 和 realized output。

\paragraph{本节小结。}
当前最可信的说法是：数据可以约束 composite route wedge、
composite innovation-arrival scale 和可能的 attenuation 强弱；
不能声称已经把所有 primitive 参数一一识别出来。

\section{Jia et al.\ 2023、Barwick et al.\ 2026 与本文机制的区分}

\subsection{为什么必须区分机制}

如果不区分机制，读者很容易把任何 post-MAH 的原研增长都解释成
“市场需求变大了”或者“审批更快了”。但这会把不同文献讲的经济对象混在一起。

\subsection{Jia et al.\ 2023：approval-delay / regulatory-capacity 机制}

Jia、Ma、Yang、Zhang（2023）强调的是 regulation for innovation，
可把它理解成更广义的 approval-delay、regulatory-capacity 或 approval-side
friction 变化。在我们的符号里，如果一个政策主要作用于审批速度或普遍审批成功，
它更接近以下对象的变化：

\begin{enumerate}[leftmargin=2em]
    \item 所有路线共享的基础成功部分 $s_i$；
    \item 或者以时间折现形式进入的等待成本；
    \item 或者广义的 approval probability，而不是特定 external route 的开放。
\end{enumerate}

这种机制的一个重要特征是：

\begin{quote}
它不必然预测 holder--producer split 上升，因为它未必改变 retained
external route 的相对吸引力。
\end{quote}

\subsection{Barwick et al.\ 2026：demand / market-reward 机制}

Barwick、Xia、Xia（2026）更接近 demand-side 或 market-reward-side 的解释。
在我们的符号里，这主要体现为
\[
R_i^I,\ R_i^E,\ R_i^T
\]
这类商业化回报对象的整体上升。若 demand 变大，所有能成功商业化的路线
都可能更有价值。

这种机制的典型预测是：

\begin{quote}
原研活动可能更强，但不必然导致 holder--producer split 上升，
因为需求扩张不是专门朝 retained external route 倾斜的。
\end{quote}

\subsection{本文 MAH 机制的特征性预测}

本文强调的 MAH route-friction mechanism 主要作用于：

\begin{enumerate}[leftmargin=2em]
    \item route set：$E$ 从无到有；
    \item route-specific friction：$\tau_i^E(M,\eta)$ 降低；
    \item route-specific implementation：$\chi_i^E(M,\eta,p)$ 提高；
    \item single-market price feedback：$p^*(M,\eta)$ 调整。
\end{enumerate}

所以它最鲜明的预测不是“所有创新都更强”，而是：

\begin{quote}
holder--producer separation 在相关 retained original-drug records 中更常见，
且这种变化最强地出现在高 $a_i$、低 $k_i$、能接触合格 CMO 支持的企业上。
\end{quote}

\subsection{三种机制的对照表}

\begin{center}
\small
\begin{tabularx}{\textwidth}{>{\raggedright\arraybackslash}p{0.23\textwidth} >{\raggedright\arraybackslash}p{0.21\textwidth} >{\raggedright\arraybackslash}p{0.22\textwidth} X}
\toprule
机制 & 在模型里主要动什么 & 对 holder--producer split 的预测 & 对原研结果的预测 \\
\midrule
Jia et al.\ 2023 式 approval-delay /
regulatory-capacity 机制 &
$s_i$、等待成本或普遍 approval friction &
不必然上升 &
可广泛提高项目推进和最终批准 \\
Barwick et al.\ 2026 式 demand 机制 &
$R_i^r$ 这类市场回报对象 &
不必然上升 &
若市场回报更大，原研激励更强 \\
MAH 路线机制 &
$\Rset_i(M)$、$\tau_i^E(M,\eta)$、
$\chi_i^E(M,\eta,p)$、$p^*(M,\eta)$ &
更容易上升，尤其在相关 retained sample 中 &
对 high-$a_i$、low-$k_i$、
external-route-relevant 企业更强 \\
\bottomrule
\end{tabularx}
\end{center}

\paragraph{本节小结。}
Jia et al.\ 2023 更像 approval-side 机制，Barwick et al.\ 2026 更像
demand-side 机制，而本文 MAH 模型是 route-set 与 route-friction 机制。
三者可能同时存在，但识别含义不同。

\section{局限、常见误解与如何正确阅读模型}

\subsection{局限}

\begin{enumerate}[leftmargin=2em]
    \item 这是部分均衡模型，不是完整行业动态均衡模型。
    \item 只闭合 CMO 支持市场，不闭合产品市场、研发投入市场和自由进入。
    \item $\zeta_i^r=s_i\chi_i^r$ 仍是复合对象，不是可被当前数据完全拆解的
    primitive system。
    \item 若缺少好的 CMO capacity proxy，很难精确识别价格反馈有多强。
    \item 若没有 application/status panel，很难把 latent opportunity arrival
    与 observed realized outcome 精细拆开。
\end{enumerate}

\subsection{常见误解一：把 MAH 当成一般 demand shock}

这是错误的。一般 demand shock 主要推动的是 $R_i^r$ 这类市场回报对象；
而 MAH 模型的关键，是新增 retained external route 并降低其特定摩擦。
如果数据里 holder--producer split 上升特别明显，这更像 route-friction
机制，而不是纯 demand 机制。

\subsection{常见误解二：把 observed share 当成简单 route probability}

这也是错误的。观测 share 必须对应 realized-flow share，因此要用
$a_i x_i^*P_i(r)\zeta_i^r$ 加权，而不是只看 $P_i(r)$。

\subsection{常见误解三：在 CMO 规划需求里乘上 \texorpdfstring{$\zeta_i^E$}{实现概率}}

这会 double count。因为 CMO 支持需求发生在路线规划阶段，
此时企业是在为可能的 retained external route 做准备。
如果再乘上 $\zeta_i^E$，就把“后续是否最终 realized”
错误地提前塞进了规划需求。

\subsection{常见误解四：把 pre-MAH 的零份额直接塞进 log odds}

如果 pre-MAH 根本没有可比 external route，零份额不是一个普通小概率事件，
而是“路线集合里没有这条路”。这时应当用 route-set expansion 的解释，
或在 post-MAH 内部利用 implementation intensity 的差异。

\subsection{常见误解五：把当前校准说成 full structural identification}

当前更诚实的说法只能是 mechanism calibration：数据约束复合对象，
而不是把每个法律成本、监督成本、实现概率、支持价格都单独拆出来。

\paragraph{本节小结。}
正确阅读这套模型的方式是：把它当成一套围绕 MAH retained external route
构建的机制框架，而不是一个什么都包的行业总模型。

\section{总小结}

整套模型可以用一条简明链条概括：

\begin{equation}
M\uparrow
\Rightarrow
\Rset_i \text{ 扩张}
\Rightarrow
\eta\uparrow \Rightarrow
\tau_i^E\downarrow,\ \chi_i^E\uparrow
\Rightarrow
\Gamma_i\uparrow
\Rightarrow
x_i^*\uparrow
\Rightarrow
a_i x_i^*\uparrow.
\label{eq:main_chain}
\end{equation}

但这条链条还有两层修正：

\begin{enumerate}[leftmargin=2em]
    \item 只有当企业真的把 $E$ 路线当成 relevant option 时，
    $\partial\Gamma_i/\partial G_i^E=P_i(E)$ 才会让改革产生明显效果；
    \item 如果更强的 external-route use 推高了 $p^*$，那么价格反馈会产生
    attenuation。
\end{enumerate}

因此，最合理的经验预期不是“MAH 让所有企业都一样受益”，而是：

\begin{quote}
MAH 最可能提高那些高研发能力、低内部生产能力、能够接触合格 CMO、
而且转让 outside option 不算太强的企业的原研激励；同时，
holder--producer split 的变化应当成为这一路线机制的重要观测痕迹。
\end{quote}

\section*{参考文献（手工列出）}

\begin{enumerate}[leftmargin=2em,label={[\arabic*]}]
    \item Acemoglu, Daron, and Joshua Linn. 2004.
    ``Market Size in Innovation: Theory and Evidence from the Pharmaceutical Industry.''
    \textit{Quarterly Journal of Economics} 119(3): 1049--1090.
    \item Arora, Ashish, Andrea Fosfuri, and Alfonso Gambardella. 2001.
    \textit{Markets for Technology: The Economics of Innovation and Corporate Strategy}.
    Cambridge, MA: MIT Press.
    \item Barwick, Panle Jia, Hongyuan Xia, and Tianli Xia. 2026.
    ``From Free Rider to Innovator: The Rise of China's Drug Development.''
    NBER Working Paper No.\ 34977.
    \item Dubois, Pierre, Olivier de Mouzon, Fiona Scott Morton, and Paul Seabright. 2015.
    ``Market Size and Pharmaceutical Innovation.''
    \textit{RAND Journal of Economics} 46(4): 844--871.
    \item Gans, Joshua S., and Scott Stern. 2003.
    ``The Product Market and the Market for Ideas:
    Commercialization Strategies for Technology Entrepreneurs.''
    \textit{Research Policy} 32(2): 333--350.
    \item Jia, Ruixue, Xiao Ma, Jianan Yang, and Yiran Zhang. 2023.
    ``Improving Regulation for Innovation: Evidence from China's Pharmaceutical Industry.''
    NBER Working Paper No.\ 31976.
    \item McFadden, Daniel. 1974.
    ``Conditional Logit Analysis of Qualitative Choice Behavior.''
    In \textit{Frontiers in Econometrics}, edited by Paul Zarembka.
    New York: Academic Press.
    \item National Medical Products Administration. 2019.
    \textit{Drug Administration Law of the People's Republic of China}.
    \item National Medical Products Administration. 2022.
    \textit{Provisions for Drug Registration}.
    \item National Medical Products Administration. 2022.
    \textit{Provisions on Supervising the Implementation of the Drug Marketing Authorization Holder's Primary Responsibility for Drug Quality and Safety}.
    \item National Medical Products Administration. 2023.
    \textit{Announcement on Strengthening Supervision and Administration of Entrusted Production by Drug Marketing Authorization Holders}.
    \item Schmookler, Jacob. 1966.
    \textit{Invention and Economic Growth}. Cambridge, MA: Harvard University Press.
    \item Teece, David J. 1986.
    ``Profiting from Technological Innovation:
    Implications for Integration, Collaboration, Licensing and Public Policy.''
    \textit{Research Policy} 15(6): 285--305.
    \item The General Office of the State Council of the People's Republic of China. 2016.
    \textit{Pilot Plan for the Marketing Authorization Holder System}.
\end{enumerate}

\end{document}
